% !TEX root = ../../../main.tex

\subsection{المرحلة الثانية: إضافة مرشح التوافقية الخامسة ($n=5$)}
تمت إضافة أول فرع ترشيح رنيني \el{($RLC$)} مصمم لاستهداف التوافقية الخامسة \el{($250 Hz$)}، وتم ضبط معاييره كما يلي: $R=0.253 \Omega$، $L=0.0121 H$، $C=3.339 \times 10^{-5} F$.

\begin{figure}[H]
    \centering
    \includegraphics[width=0.8\linewidth]{04 - 1st filter parameters.png}
    \caption{نافذة بارامترات مرشح التوافقية الخامسة في ماتلاب}
\end{figure}
\begin{figure}[H]
        \centering
        \includegraphics[width=1.05\linewidth]{03 - Circuit with 1st filter.png}
        \caption{شكل الدارة بعد إضافة المرشح الأول ($n=5$)}
    \end{figure}\newpage

عند إضافة المرشح الأول ($n=5$)، نلاحظ من خلال أداة التحليل تلاشي القمم المقابلة لهذا التردد في طيف التردد:
\begin{figure}[H]
  \centering
  \includegraphics[width=0.9\textwidth]{05 - CurrentData with 1st filter.png}
    \caption{\small نتائج تحليل \el{FFT} للتيار بعد اضافة المرشح الأول - تظهر \el{THD} بنسبة $14.83\%$}
    \label{fig:05}
\end{figure}
\begin{figure}[H]
  \centering
  \includegraphics[width=0.9\textwidth]{06 - VoltageData with 1st filter.png}
    \caption{\small نتائج تحليل \el{FFT} للجهد بعد اضافة المرشح الأول - تظهر \el{THD} بنسبة $7.27\%$}
    \label{fig:06}
\end{figure}

  \newpage

\begin{itemize}
  \item أظهرت النتائج في الصورة (الشكل \ref{fig:05}) انخفاضاً ملموساً في نسبة تشوه التيار بعد اضافة المرشح الأول، لتصل إلى \el{\boldmath $14.83\%$}.
  \item أظهرت النتائج في الصورة (الشكل \ref{fig:06}) استقرار نسبة تشوه الجهد عند \el{\boldmath $7.27\%$} (بناءً على بارامترات ممانعة الشبكة المستخدمة في النموذج).
\end{itemize}

بعد المحاكاة، لوحظ انخفاض ملموس في نسبة التشوه التوافقي، مما يؤكد فعالية المرشح في امتصاص تيارات الرتبة الخامسة.
